% STATUS: rewritten 2026-10-12 (owner: the figure showed too few kinds; list them, and let each chapter add its related ones)
% Sources read in full: Hochschild 1942 (p. 688), Weibel 1994 (Definition 6.4.1), Benkart, Lopes and Ondrus 2015 (p. 47). Every other relative is defined in the chapter named below, with its sources.
\section{The kinds of derivation-like maps}\label{sec:basic-theory:family-tree}

We have so far studied one rule, the product rule. We now collect the ways in
which later authors changed it.

Some changes keep the rule and move the target. Hochschild allowed the values of
$D$ to lie in a larger algebra \cite[Definition~3.1, p.~688]{Hochschild1942}.
Weibel replaced the algebra by a group, and the product rule by
$D(gh) = gD(h) + D(g)$ \cite[Definition~6.4.1]{Weibel1994}. We met both changes
in Section~\ref{sec:basic-theory:inner-outer}.

We turn to changes that keep the shape of the rule and twist one factor.
Section~\ref{sec:origins:ore-and-twists} met maps $\delta$ with
$\delta(rs) = \delta(r)s + \sigma(r)\delta(s)$, where $\sigma$ is a ring map from
$R$ to itself \cite[p.~47]{BenkartLopesOndrus2015}. Chapter~\ref{ch:generalized-twisted}
treats these maps and their generalizations.

We cannot show the kinds in one picture, because there are too many and the list grows with the literature. We list them by the change they make to the rule, and we name the chapter that treats each one. Every chapter defines its kinds from the sources, gives the
known implications and quotes the counterexamples.

\begin{itemize}
\item \emph{A new target.} Derivations into a bimodule or a larger algebra, and
crossed homomorphisms (Section~\ref{sec:basic-theory:inner-outer}); derivations
of groups (Chapter~\ref{ch:cohomology}); the universal derivation and K\"ahler
differentials (Chapter~\ref{ch:commutative}).
\item \emph{A twisted factor.} Skew, $\sigma$-, $\alpha$- and
$(\sigma,\tau)$-derivations, and generalized derivations
(Chapter~\ref{ch:generalized-twisted}); semiderivations, left derivations,
reverse derivations and multiplicative derivations
(Chapter~\ref{ch:generalized-twisted}); $\delta$-derivations, quasi-derivations
and anti-derivations (Chapter~\ref{ch:nonassociative}).
\item \emph{A new product.} Jordan derivations, Jordan triple derivations,
Jordan $n$-derivations, Jordan $*$-derivations and Jordan left derivations
(Chapter~\ref{ch:jordan-type}); Lie derivations, Lie triple derivations, Lie
$n$-derivations and the nonlinear Lie-type maps (Chapter~\ref{ch:lie-type}).
\item \emph{Maps built from derivations.} Higher derivations and
Hasse--Schmidt sequences, biderivations, triderivations, $n$-derivations and
multiderivations (Chapter~\ref{ch:higher-multi}); exponential maps and locally nilpotent derivations (Section~\ref{sec:basic-theory:exponential} and
Chapter~\ref{ch:commutative}).
\item \emph{A condition at each point.} Local derivations, $2$-local
derivations and derivable maps (Chapter~\ref{ch:local-characterizations});
commuting and centralizing maps and functional identities
(Chapter~\ref{ch:functional-identities}).
\item \emph{Kinds named in the titles of the collected literature.} The titles
of the papers gathered for this book name further maps, each in at least
six titles. Banach and operator algebras
(Chapter~\ref{ch:banach-operator}): closed, unbounded, normal, point and
double derivations, superderivations, and Hom-derivations. Nonassociative
algebras (Chapter~\ref{ch:nonassociative}): ternary derivations. Lie algebras
(Chapter~\ref{ch:lie-algebras}): prederivations and central derivations.
Commutative algebra (Chapter~\ref{ch:commutative}): integrable and quadratic
derivations. Matrix and triangular algebras
(Chapter~\ref{ch:matrix-triangular}): additive derivations of nest algebras.
Functional identities (Chapter~\ref{ch:functional-identities}): commuting and
hypercommuting derivations. Polynomial and Weyl algebras
(Chapter~\ref{ch:polynomial-weyl}): right derivations. Partial derivations of
functions of $n$ variables. Each chapter defines its kinds from the sources it cites.
\item \emph{Maps indexed by other maps.} $f$-derivations and
$(f,g)$-derivations, with their permuting, symmetric and tri-$f$ forms, in
lattices, BCI-algebras and subtraction algebras (Chapter~\ref{ch:other-systems}),
and $f$-derivations on rings and modules (Chapter~\ref{ch:generalized-twisted});
$(\alpha,\beta)$-, $(\theta,\varphi)$- and $(\sigma,\tau)$-derivations in their
Jordan, Lie and generalized forms (Chapters~\ref{ch:jordan-type},
\ref{ch:lie-type} and~\ref{ch:generalized-twisted}). Derivations of almost
$f$-algebras (Chapter~\ref{ch:banach-operator}) take their name from a different
$f$: it refers to a Riesz space, not to a map.
\item \emph{Further names from the collected titles.} Each occurs in at least six
titles. Order bounded, positive, simple, linear, free, ring, almost inner,
generalized inner, orthogonal generalized, locally finite and local triple
derivations; Jordan and Lie higher derivations; generalized Jordan left
derivations. Each chapter defines those of its setting from the sources it cites.
\item \emph{More names from the collected titles.} A second pass over the
titles finds names in at least three titles each. Chapter~\ref{ch:generalized-twisted}:
$b$-generalized, $\alpha$-generalized, $q$-skew, generalized semiderivations,
bi-semiderivations and multiplicative semiderivations.
Chapter~\ref{ch:prime-semiprime}: $(\alpha,\alpha)$-, $(\theta,\theta)$-,
$(\alpha,\beta,\gamma)$- and $(\sigma,\tau,\xi)$-derivations,
$(\alpha,\beta)$-reverse derivations, $g$-derivations, $3$-derivations and left
biderivations. Chapter~\ref{ch:higher-multi}: symmetric biderivations,
permuting triderivations, $(m,n)$-derivations and $k$-derivations.
Chapter~\ref{ch:commutative}: monomial, Weitzenb\"ock, elementary, semisimple,
Jacobian, logarithmic, factorizable and $r$-derivations.
Chapter~\ref{ch:nonassociative}: conformal, homogeneous, special and
$(n+1)$-ary derivations, and Lie and Jordan superderivations.
Chapter~\ref{ch:banach-operator}: continuous, compact, random, hyponormal,
self-adjoint and quasi-free derivations. Chapter~\ref{ch:lie-algebras}:
nonsingular, symplectic and locally finite derivations.
Chapter~\ref{ch:other-systems}: $p$-derivations and $\gamma$-derivations.
Chapter~\ref{ch:functional-identities}: two-sided centralizers, super-commuting
maps and $n$-commuting maps. Chapter~\ref{ch:cohomology}: $2$-cocycles. The
file lists every name with its count. Each chapter defines its kinds from the
sources it cites.
\item \emph{A different algebraic setting.} Prime and semiprime rings
(Chapter~\ref{ch:prime-semiprime}). Lie algebras
(Chapter~\ref{ch:lie-algebras}). Jordan, alternative and Leibniz algebras,
superalgebras, Poisson algebras, Hom-Lie algebras and $n$-ary algebras
(Chapter~\ref{ch:nonassociative}). Hopf algebras, with coderivations
(Chapter~\ref{ch:group-hopf}). Differential and difference rings
(Chapter~\ref{ch:differential-algebra}). Banach and operator algebras, with
approximate derivations (Chapter~\ref{ch:banach-operator}). Near-rings, semirings,
Gamma-rings and lattices (Chapter~\ref{ch:other-systems}).
\end{itemize}

Section~\ref{sec:basic-theory:question-families} makes precise the one kind of
question that the list raises: when does one kind force another?
